The rapid growth of digital education has created a need for simple tools that help instructors build interactive learning content without specialist technical skills. At Vietnam National University -- Ho Chi Minh City (VNU-HCM), creating H5P interactive videos through the WordPress ecosystem remains cumbersome and slow for many educators.

This thesis presents the design and evaluation of a dedicated platform for AI-assisted H5P video authoring. The system combines a React~18 and TypeScript frontend, a Node.js/Express backend, a SQLite database, and a transcript-driven AI pipeline that generates candidate H5P interactions for instructor review. The pipeline accepts caption sources from uploaded files, YouTube captions, or audio transcription, then validates each suggestion before it reaches the editor.

A comparative usability study with ten Information Technology students showed that the platform substantially improved authoring efficiency and usability over the WordPress baseline. Task completion increased from 30\% to 95\%, median creation time dropped from 45 minutes to 3.2 minutes, and usability scores rose from 2.1 to 4.8 on a five-point scale. The results indicate that combining user-centred design with generative AI can lower the barrier to creating interactive pedagogical material for university teaching.

\bigskip
\noindent
\textbf{Keywords:} H5P, interactive video, educational technology, AI-assisted content authoring, usability evaluation, user-centred design.
